% page 136 of the facsimile (image p-135 of the scan)
% block 160.0 x 237.7 mm ; left margin 8.0 mm
\pgc{-12.700mm}{0.000mm}{
\l{\cel{58}126}
\l{}
\l{}
\l{\sou{edukar}.- (\sou{trans}.) I. Formacar infanto, puero, adolecanto,}
\l{\cel{3}per developar e direktar lua fakultati korpala, morala}
\l{\cel{3}(etikala), mentala, religiala. - II. Kustumigar ulu pri}
\l{\cel{3}la ceremonialo mondumala (politeso, konveni, deci, e c.)}
\l{\cel{3}- III. Exercar la animali pri lia kapablesi inata. -}
\l{\cel{3}IV. Nutrar, plumultigar, entratenar ula animali por lia}
\l{\cel{3}manjeso, lia laboro e la plezuro quan onu recevas de li ;}
\l{\cel{3}same pri planti. - EFIS.}
\l{}
\l{\sou{efacar}. - (\sou{trans}.)Desaparigar to quo trasesis. - EF.}
\l{}
\l{\sou{efekto.}\cel{1}- Rezultajo de la efiko di kauzo. - DEFIRS.}
\l{}
\l{\sou{efektivo}. I. (\sou{armeo, administrantaro}). La nombro de la}
\l{\cel{3}individui qui fakte servas. - II. La nombro de la indi-}
\l{\cel{3}vidui di qui la nomi enrejistresis en la etati di societo,}
\l{\cel{3}e c. - DEFIS.}
\l{}
\l{\sou{efemera}.- Qua duras dum nur un dio. - DEFIS.}
\l{}
\l{\sou{efemerido}.- I. Tabeli astronomiala qui indikas, pri singla}
\l{\cel{3}dio di la yaro, la situeso di la astri. - II. Libro qua}
\l{\cel{3}indikas anticipe la fakti astronomiala o meteorologiala,}
\l{\cel{3}kalkulebla e predicebla pri ta o ca epoko. - DEFIS.}
\l{}
\l{\sou{efervecar}. (\sou{netrans}.) I. (\sou{aludante gaso}). Ebuliar tra liquido.}
\l{\cel{3}II. Emocar vivace, ma neduronte. - EFIS.}
\l{}
\l{\sou{eficienta}.- Qua produktas efekto maxima po laboro, esforco,}
\l{\cel{3}energio, spenso, minima. - EF.}
\l{}
\l{\sou{efikar}. - (\sou{netrans}.)Produktar la efekto quan onu expektas de}
\l{\cel{3}lu. - EFIS.}
\l{}
\l{\sou{eflorecar}.- (\sou{netrans.})\cel{2}I. (\sou{kemio}) La efekto qua konsistas ye}
\l{\cel{3}la formaco di kristali ye la surfaco di solido o liquido.}
\l{\cel{3}- II. (\sou{fiziol.}) Fenomeno qua karakterizesas da mikra infl-}
\l{\cel{3}uro di la pelo. - DEFIS.}
\l{}
\l{\sou{-eg-}\cel{1}.\cel{2}Sufixo qua indikas \textquotedbl{}la grado maxim alta\textquotedbl{}.}
\l{}
\l{\sou{egala}.\cel{2}Same-quanta o same-quala kam to a quo onu komparas}
\l{\cel{3}lu. - DEFIRS.}
\l{}
\l{\sou{egarir}.(trans.) Lokizir distrakte (ulo) ube onu ne plus tro-}
\l{\cel{3}vas lu.-F.}
\l{}
\l{\sou{egardar.(trans}.) Atencar kozo o persono, pro lua importo. - EF.}
\l{}
\l{\sou{egido}.- I. (en\cel{1}\sou{la epoki antiqua}). Shildo di Pallas, kovrita}
\l{\cel{3}per la felo di la kaprino Amalte, e sur qua reprezentesabis}
\l{\cel{3}la kapo di Meduzo. - II. To quo protektas. - DEFIS.}
}
